% Part: lambda-calculus
% Chapter: church-rosser
% Section: definitions-and-properties

\documentclass[../../../include/open-logic-section]{subfiles}

\begin{document}

\olfileid{lam}{cr}{dap}

\olsection{تعريف او خاصيتونه}

په دې څپرکي کښې د چرچ--روسر خاصيت او د هغه ځينې عامې
ځانګړنې معرفي کوو۔

\begin{defn}[د چرچ--روسر خاصيت، CR] پر ترمونو يوه اړيکه
  $\xredone$ هله او يوازې هله د \emph{چرچ--روسر خاصيت} لري
  چې کله هم $M \xredone P$ او $M \xredone Q$ وي، داسې يو~$N$
  موجود وي چې $P \xredone N$ او $Q \xredone N$ وي۔
\end{defn}

لامبډا حساب د محاسبې د يوه ماډل په توګه ليدلے شو: عادي بڼې
لرونکي ترمونه يې «ارزښتونه» دي او پر ترمونو د راکمېدو اړيکه يې
«د محاسبې قاعدې» دي۔ د چرچ--روسر خاصيت وايي چې که محاسبه په
څو لارو پرمخ تللے شي، نو د هغې عادي پايله، که موجوده وي، يوازې
يوه وي۔ \textbf{د سرچينې سپيناوی (OLLAM-035):} دا خاصيت د
عادي پايلې موجوديت نۀ تضمينوي؛ د يکتايي خبره د موجوديت په شرط ده۔

د ابتدايي الجبره د بېلګې په توګه، $4 \times (1+2) + 3$ په څو
لارو حسابېدے شي۔ يا ئې $4 \times 3+3$ ته راکمولاے شو، که
لومړے $1+2$،~$3$ ته راکم کړو؛ يا ئې $4 \times 1+4
\times 2+3$ ته راکمولاے شو، که لومړے پر $4 \times (1+2)$
د وېش قانون وکاروو۔ خو دواړه لارې بيا $12+3$ ته راکمېږي۔

که $\xredone$ د $\beta$-راکمول اړيکه ونيسو، په اسانه وينو
چې د چرچ--روسر د خاصيت يوه پايله دا ده: که يو ترم عادي بڼه
ولري، هغه يکتا ده۔ فرض کړئ $M$، $P$ او $Q$ ته راکمېږي
او دواړه عادي بڼې دي۔ د چرچ--روسر د خاصيت له مخې داسې يو
$N$ شته چې $P$ او $Q$ دواړه ورته راکمېږي۔ څرنګه چې د فرض
له مخې $P$ او $Q$ عادي بڼې دي، نو له $P$ او $Q$ څخه $N$
ته راکمول يوازې بې‌بدلونه راکمول کېدے شي؛ يعنې $P$، $Q$
او $N$ يو شان دي۔ ځکه د يوه ترم د \emph{هماغې} عادي بڼې
په اړه خبرې کول سم دي۔

که لامبډا حساب د محاسبې ماډل وګڼو، د يوه ترم عادي بڼه د
هغه ترم له محاسبې څخه د راوتلې «وروستۍ پايلې» په توګه ليدلے
شو۔ د پورته پايلې معنا دا ده چې که وروستۍ پايله وي، نو يوازې
يوه وي؛ لکه د $4 \times (1+2)+3$ د حساب پايله چې~$15$ ده۔

\begin{thm} \ollabel{thm:str}
  که يوه اړيکه $\xredone$ د چرچ--روسر خاصيت ولري او $\xred$
  تر ټولو کوچنۍ انتقالي اړيکه وي چې $\xredone$ پکې شامله ده،
  نو $\xred$ هم د چرچ--روسر خاصيت لري۔
\end{thm}

\begin{proof}
  فرض کړئ
  \begin{align*}
    M & \xredone P_1 \xredone \dots \xredone P_m \text{ او}\\
    M & \xredone Q_1 \xredone \dots \xredone Q_n.
  \end{align*}
  قضيه د ترمونو د يوه $N$ جال په جوړولو ثابتوو چې لوړوالے يې
  $m + 1$ او پلنوالے يې $n + 1$ دے۔ د $i$-مې ليکې او
  $j$-مې ستنې ترم په $N_{i,j}$ ښيو۔
  
  $N$ داسې جوړوو چې $N_{i,j} \xredone N_{i+1,j}$ او
  $N_{i,j} \xredone N_{i,j+1}$ وي۔ تعريف يې دا دے:
  \begin{align*}
    N_{0,0} &= M \\
    N_{i,0} &= P_i && \text{که } 1 \le i \le m \\
    N_{0,j} &= Q_j && \text{که } 1 \le j \le n \\
  \intertext{او په نورو حالاتو کښې:}
    N_{i,j} &= R
  \end{align*}
  دلته $R$ داسې ترم دے چې $N_{i-1,j} \xredone R$ او
  $N_{i,j-1} \xredone R$ وي۔ د $\xredone$ د چرچ--روسر د
  خاصيت له مخې داسې ترم تل موجود وي۔

  اوس $N_{m,0} \xredone \dots \xredone N_{m,n}$ او $N_{0,n}
  \xredone \dots \xredone N_{m,n}$ لرو۔ پام وکړئ چې $N_{m,0}$
  هماغه $P$ او $N_{0,n}$ هماغه $Q$ دے۔ د $\xred$ له تعريفه
  قضيه ثابتېږي۔ \textbf{د سرچينې سپيناوی (OLLAM-036):}
  په ثبوت کښې دا دوه توري د پورته دوو لړيو وروستي ترمونه
  نوموي؛ سرچينه يې تر دې جملې وړاندې په څرګنده نۀ پېژني۔
\end{proof}

\end{document}
